\chapter{Происхождение инвариантного родительского функционала логарифма определителя из меры}
\label{ch:v9-endpoint-kms-logdet-measure-origin}

\section{Проверяемый вопрос}

Предыдущий гейт построил два положительных type-operator
\begin{equation}
 R_\theta=\operatorname{diag}(r_{\theta s},r_{\theta a},
 r_{\theta t},r_{\theta t},r_{\theta t}),
 \qquad
 R_\kappa=\operatorname{diag}(r_{\kappa s},r_{\kappa a},
 r_{\kappa t},r_{\kappa t},r_{\kappa t})
 \label{eq:v9-kms-logdet-measure-type-operators}
\end{equation}
и условный parent
\begin{equation}
 B_{\rm target}=-\log\det R_\theta-\log\det R_\kappa.
 \label{eq:v9-kms-logdet-measure-target}
\end{equation}
Нужно проверить, является ли этот term Jacobian или effective action уже
имеющейся меры, а не новым постулатом.

\section{Gaussian sign audit}

Для positive Hermitian $R$ комплексная bosonic Gaussian integral имеет
determinant в знаменателе:
\begin{equation}
 Z_{\rm B}(R)\propto(\det R)^{-1},
 \qquad
 -\log Z_{\rm B}(R)=+\log\det R+\operatorname{const}.
 \label{eq:v9-kms-logdet-measure-complex-boson}
\end{equation}
Знак противоположен требуемому. Real boson даёт половину того же знака:
\begin{equation}
 Z_{\rm RB}(R)\propto(\det R)^{-1/2},
 \qquad
 -\log Z_{\rm RB}(R)=+\frac12\log\det R+\operatorname{const}.
 \label{eq:v9-kms-logdet-measure-real-boson}
\end{equation}

Для пяти complex Grassmann-пар формула Березина даёт
\begin{equation}
 Z_{\rm F}(R)=
 \int D\bar\psi\,D\psi\,
 \exp\!\left(-\bar\psi R\psi\right)=\det R.
 \label{eq:v9-kms-logdet-measure-berezin}
\end{equation}
Поэтому effective action имеет в точности нужный знак и коэффициент:
\begin{equation}
 S_{\rm F,eff}(R)=-\log Z_{\rm F}(R)=-\log\det R.
 \label{eq:v9-kms-logdet-measure-fermion-sign}
\end{equation}
Majorana/Pfaffian-вариант даёт только половинный coefficient,
\begin{equation}
 \operatorname{Pf}(A)^2=\det A,
 \qquad
 -\log\operatorname{Pf}(A)=-\frac12\log\det A,
 \label{eq:v9-kms-logdet-measure-pfaffian}
\end{equation}
и без удвоения не совпадает с target.

\section{Минимальная auxiliary размерность}

Multiplicity pattern $1\oplus1\oplus3$ фиксирует determinant
\begin{equation}
 \det R=r_s r_a r_t^3,
 \qquad
 \deg(\det R)=5.
 \label{eq:v9-kms-logdet-measure-degree-five}
\end{equation}
Следовательно, минимальный linear Gaussian carrier одного package содержит
пять complex Grassmann-пар. Для двух независимых packages нужен block
\begin{equation}
 D_{\rm aux}=R_\theta\oplus R_\kappa,
 \qquad
 \dim_{\mathbb C}D_{\rm aux}=10.
 \label{eq:v9-kms-logdet-measure-aux-block}
\end{equation}
Его determinant факторизуется:
\begin{equation}
 \det D_{\rm aux}=\det R_\theta\,\det R_\kappa,
 \qquad
 -\log\det D_{\rm aux}=B_{\rm target}.
 \label{eq:v9-kms-logdet-measure-block-factorization}
\end{equation}
Одна копия каждого блока даёт coefficient one; $n_\theta,n_\kappa$
копий дали бы
\begin{equation}
 S_{\rm eff}=-n_\theta\log\det R_\theta
 -n_\kappa\log\det R_\kappa,
 \qquad n_\theta,n_\kappa\in\mathbb N.
 \label{eq:v9-kms-logdet-measure-copy-number}
\end{equation}

\section{No-go обычного coordinate Jacobian}

В log-ratio chart
\begin{equation}
 (r_s,r_a,r_t)=
 \frac5{e^u+e^v+3}(e^u,e^v,1)
 \label{eq:v9-kms-logdet-measure-chart}
\end{equation}
двумерный Jacobian равен
\begin{equation}
 J_{\rm chart}
 =\det\frac{\partial(r_s,r_a)}{\partial(u,v)}
 =\frac35 r_s r_a r_t.
 \label{eq:v9-kms-logdet-measure-chart-jacobian}
\end{equation}
Target determinant содержит triplet multiplicity three:
\begin{equation}
 \frac{\det R}{J_{\rm chart}}
 =\frac{r_s r_a r_t^3}{(3/5)r_s r_a r_t}
 =\frac53 r_t^2.
 \label{eq:v9-kms-logdet-measure-chart-mismatch}
\end{equation}
Отношение непостоянно. Поэтому flat measure в $(u,v)$ или индуцированная
coordinate measure не порождает invariant barrier.

\section{Carrier audit}

Существующая creation-cell имеет только physical source и пять channel
targets:
\begin{equation}
 \mathcal H_{\rm cr}
 =\mathbb C|0\rangle\oplus
 \mathbb C|s\rangle\oplus\mathbb C|a\rangle\oplus\mathbb C^3|t\rangle,
 \qquad \dim_{\mathbb C}\mathcal H_{\rm cr}=6.
 \label{eq:v9-kms-logdet-measure-creation-cell}
\end{equation}
В её зарегистрированных полях нет независимого odd auxiliary module,
Berezin measure или gauge map с Faddeev--Popov Jacobian $R$. Значит,
найденная fermionic representation является conditional extension:
\begin{equation}
 N_{\rm conditional\ measure\ mechanism}=1/1,
 \qquad
 N_{\rm inherited\ auxiliary\ measure}=0/1.
 \label{eq:v9-kms-logdet-measure-carrier-ledger}
\end{equation}

\section{ProofDSL certificate}

LCF-ядро проверяет determinant identities, homogeneous degrees, exact
effective signs и coordinate-Jacobian mismatch:
\begin{equation}
 N_{\rm ProofDSL\ obligations}=10/10,
 \qquad
 N_{\rm registered\ gates}=50,
 \qquad
 N_{\rm registered\ obligations}=382.
 \label{eq:v9-kms-logdet-measure-proofdsl}
\end{equation}
Gaussian и Berezin integration identities остаются внешними аналитическими
леммами; kernel сертифицирует их конечномерные determinant-следствия.

\section{Итог}

\begin{equation}
 \begin{aligned}
 N_{\rm measure\ candidates}&=1/7,\\
 N_{\rm conditional\ fermionic\ representation}&=1/1,\\
 N_{\rm minimal\ auxiliary\ dimension}&=10,\\
 N_{\rm inherited\ logdet\ origin}&=0/1,\\
 N_{\rm physical\ four\ slot\ parent}&=0/1.
 \end{aligned}
 \label{eq:v9-kms-logdet-measure-final-ledger}
\end{equation}

\begin{theorem}[Условная measure-реализация логарифм определителя parent]
Две independent complex fermionic Gaussian measures на блоках размерности
$5+5$ порождают ровно
$-\log\det R_\theta-\log\det R_\kappa$. Степень determinant показывает
минимальность десяти complex Grassmann-пар для linear block realization.
\end{theorem}

\begin{nogocriterion}[Механизм не равен происхождению поля]
Ни flat/coordinate measure, ни bosonic Gaussian не порождают target term.
Fermionic determinant даёт правильную формулу, но требуемый odd auxiliary
module отсутствует в текущем four-slot carrier. Поэтому physical origin
логарифм определителя parent остаётся $0/1$.
\end{nogocriterion}

\begin{protocol}\begin{enumerate}
\item target: \textbf{$-\log\det R_\theta-\log\det R_\kappa$};
\item correct sign: \textbf{complex fermionic Gaussian};
\item minimal carrier: \textbf{$5+5=10$ complex Grassmann-пар};
\item coordinate Jacobian: \textbf{$(3/5)r_sr_ar_t\neq r_sr_ar_t^3$};
\item ProofDSL: \textbf{$10/10$, registry $50/382$};
\item inherited measure-origin: \textbf{$0/1$};
\item следующий гейт: \textbf{auxiliary fermion module admission}.
\end{enumerate}\end{protocol}

Машинный сертификат сохранён в
\path{s2t_v9_endpoint_creation_kms_relative_shape_logdet_parent_measure_origin_gate_results.json}.

\begin{thebibliography}{9}
\bibitem{Berezin1966}
F.~A.~Berezin,
\emph{The Method of Second Quantization},
Academic Press, 1966.
\end{thebibliography}